% !TeX root = ../thuthesis-example.tex


\chapter{本文编写的核心程序代码}
\begin{lstlisting}[caption={交互式相机标定流程},
label={lst:camera_calibration}]

Algorithm: Interactive Camera Calibration Pipeline

Input:
camera index c
chessboard pattern size (W,H)

Output:
intrinsic matrix K
distortion coefficients d

Initialize GUI interface
Create image storage directory

Set system states:
camera_opened <- False
undistort_mode <- False
K <- empty
d <- empty

while application is running do

```
if user requests "Open Camera" then
    Open camera stream using c
    Start frame acquisition thread
    Enable capture and calibration controls
end if

if camera stream available then
    Acquire newest frame

    if undistort_mode = True
    and calibration parameters exist then
        Apply image undistortion
    end if

    Render frame to GUI
end if

if user requests image capture then
    Save current frame
    Refresh captured image list
end if

if user selects an image then
    Load selected image
    Display image preview
end if

if user requests chessboard detection then
    Load latest captured image
    Convert image to grayscale
    Detect chessboard corners using (W,H)

    if corners detected then
        Refine corners to subpixel precision
        Draw chessboard visualization
        Save visualization result
    else
        Report detection failure
    end if
end if

if user requests camera calibration then

    Construct chessboard object points

    Initialize:
        O <- empty object points
        I <- empty image points

    for each captured image do
        Convert image to grayscale
        Detect chessboard corners

        if corners detected then
            Refine corner positions
            Append object points to O
            Append image points to I
        end if
    end for

    if number of valid images < 3 then
        Report insufficient calibration data
    else
        (K,d) <- CalibrateCamera(O,I)
        Compute reprojection error
        Display calibration quality metrics
    end if
end if

if user requests parameter saving then
    Save K, d, and image size
end if

if user requests parameter loading then
    Load calibration parameters
    Enable undistortion preview
end if

if user toggles undistortion preview then
    Toggle undistort_mode
end if
```

end while

Release camera resource
Terminate application

\end{lstlisting}


\begin{lstlisting}[caption={笔记本电脑接收图传、发送控制信号核心代码},
label={lst:ground_station_control}]

Algorithm: Ground Control Interface for CRSF-Based UAV Teleoperation

Input:
serial port P
baudrate B
camera index c

Output:
continuous CRSF command transmission
visual feedback for teleoperation

Initialize GUI interface
Initialize serial communication (P,B)
Open camera stream using c

Initialize slider states:
(roll, pitch, throttle, yaw)
<- (0, 0, -1, 0)

Initialize system states:
armed <- False
arm_request <- False
angle_mode <- False

while system is running do

```
Poll keyboard and mouse events

if camera switch key pressed then
    Release current camera stream
    Open selected camera
end if

if user clicks ARM button then
    if armed = True then
        arm_request <- False
    else if throttle is sufficiently low then
        arm_request <- True
    end if
end if

if user clicks ANGLE button then
    Toggle angle_mode
end if

for each control slider do
    if slider is dragged then
        Update normalized control value
        Clamp value into interval [-1,1]
    end if
end for

for roll, pitch, yaw channels do
    if slider is not dragged then
        Gradually decay command toward zero
    end if
end for

Convert normalized commands into CRSF channels:
    ch[0:3] <- MapToCRSF(
                    roll,
                    pitch,
                    throttle,
                    yaw)

Determine throttle safety state:
    throttle_low <- (ch[2] <= T_arm)

if arm_request = True
and throttle_low = True then
    armed <- True
else if arm_request = False then
    armed <- False
end if

Update auxiliary switch channels:

    ch[5] <- MID if armed
             else LOW

    ch[6] <- MID if angle_mode
             else LOW

if transmission period elapsed then
    Pack RC channels into CRSF payload
    Compute CRC checksum
    Send CRSF frame through serial port
end if

if camera frame available then
    Acquire image frame
    Resize and convert image format
    Render image onto GUI
end if

Render sliders, buttons,
channel values and status messages

Refresh GUI display
```

end while

Release camera resource
Release serial resource
Terminate application

\end{lstlisting}
